\documentclass[11pt,a4paper]{article}
\usepackage{kotex}
\usepackage{hyperref}
\hypersetup{hidelinks}
\begin{document}
\section{SUMMARY}
미확인 사항을 보존한다.
\section{분석 배경 및 KV cache 문제}
미확인 사항을 보존한다.
\section{기술 선정 및 선정 이유}
미확인 사항을 보존한다.
\section{RDKV·Photonic-CXL 기술 개요}
미확인 사항을 보존한다.
\subsection{RDKV}
미확인 사항을 보존한다.
\subsection{Photonic-CXL}
미확인 사항을 보존한다.
\section{평가 기준 및 분석 방법}
미확인 사항을 보존한다.
\section{관점별 평가}
미확인 사항을 보존한다. RDKV\cite{SW01_RDKV}, Photonic-CXL\cite{HW01_PHOTONIC_CXL}.
\subsection{기술 성숙도}
공개 정보 기반 팀 추정이며 공식 인증이 아니다.
% BEGIN_TRL_ASSESSMENT SW-01
\paragraph{RDKV}
추정 TRL: 4. GPU 실험의 구성요소 검증을 근거로 한다.\cite{SW01_RDKV}
다음 미확인 조건: 대표 QA 워크로드의 요구 성능 검증.
목표 서비스의 동시성 조건 미확인.
% END_TRL_ASSESSMENT SW-01
% BEGIN_TRL_ASSESSMENT HW-01
\paragraph{Photonic-CXL}
추정 TRL: 미확인. 단계 판정 자료가 충분하지 않아 숫자를 확정하지 않는다.
다음 미확인 조건: 구현 버전의 단계별 원문 근거.
제공 자료에서 단계 판정 근거 미확인.
% END_TRL_ASSESSMENT HW-01

\subsection{시장성}
미확인 사항을 보존한다.
\subsection{이해관계자}
미확인 사항을 보존한다.
\subsection{도메인 적용성}
미확인 사항을 보존한다.
\section{관점 간 비교 및 시사점}
미확인 사항을 보존한다.
\subsection{일치하는 평가}
미확인 사항을 보존한다.
\subsection{상충하는 평가}
미확인 사항을 보존한다.
\subsection{조건별 상대적 적합성}
미확인 사항을 보존한다.
\subsection{병행 가능성}
미확인 사항을 보존한다.
\section{한계점 및 확증편향 방지}
미확인 사항을 보존한다. 모의 상위 Agent 입력이며 실제 성능 재현 결과가 아니다.
\subsection{현재 자료의 한계}
미확인 사항을 보존한다.
\subsection{도입 판단 전 검증 우선순위}
미확인 사항을 보존한다.
\section{REFERENCE}

\renewcommand{\refname}{}
\begin{thebibliography}{9}
\bibitem{SW01_RDKV} RDKV paper.
\bibitem{HW01_PHOTONIC_CXL} Photonic-CXL paper.
\end{thebibliography}

\end{document}
